\documentclass[11pt]{article}
\usepackage{cs1200}
\usepackage[normalem]{ulem}

\begin{document}

\psHeader{9}{Tuesday April 21, 2026 (11:59pm)}

Please see the syllabus for the full collaboration and generative AI policy, as well as information on grading, late days, and revisions.

All sources of ideas, including (but not restricted to) any collaborators, AI tools, people outside of the course, websites, ARC tutors, and textbooks other than Hesterberg--Vadhan must be listed on your submitted homework along with a brief description of how they influenced your work. You need not cite core course resources, which are lectures, the Hesterberg--Vadhan textbook, sections, SREs, problem sets and solutions sets from earlier in the semester. If you use any concepts, terminology, or problem-solving approaches not covered in the course material by that point in the semester, you must describe the source of that idea. If you credit an AI tool for a particular idea, then you should also provide a primary source that corroborates it. Github Copilot and similar tools should be turned off when working on programming assignments.

If you did not have any collaborators or external resources, please write 'none.' Please remember to select pages when you submit on Gradescope. A problem set on the border between two letter grades cannot be rounded up if pages are not selected OR collaborators not noted.

\emph{Note that, historically, people have found this Pset to be one of the more difficult and time-consuming ones, so please plan accordingly (start early, come to OH and section, and ask questions on Ed!) and don't be discouraged! This is a hard topic, but lots of practice will help.}

\vspace{1em}

\textbf{Your name: }

\textbf{Collaborators and External Resources:}

\textbf{No. of late days used on previous psets: }

\textbf{No. of late days used after including this pset: }

\vspace{0.2in}

\noindent The purpose of this problem set is to to reinforce the definitions and basic theory of our complexity classes and practice $\NP$-completeness proofs and related reductions.

\begin{enumerate}
    \item (Complexity Classes and Reductions)
        Consider the computational problem $\Pi = (\Inputs,\Outputs,f)$ where $\Inputs=\Outputs=\N$ and $f(x)=\N$ for all $x\in \N$.  Show that $\Pi\in \Psearch$ but $\Pi\notin \NPsearch$.  Thus, $\Psearch\nsubseteq\NPsearch$.

    \item ($\NPsearch$-completeness) Consider the following variant of $k$-SAT.

       \compprob{\FlipkSAT}
        {A $k$-CNF formula $\varphi(x_0,\ldots,x_{n-1})$}
        {An assignment $\alpha\in\zo^n$ such that $\varphi(\alpha)=1$ and $\varphi(\neg \alpha)=1$ (if one exists), where $\neg \alpha$ is the bitwise negation of $\alpha$}

        Your proofs for this problem should all include correctness (in both directions) and runtime, following the outline of proofs done in lecture and the textbook.

        \begin{enumerate}
        \item Prove that \FlipFourSAT\ is $\NPsearch$-complete by reduction from \ThreeSAT.  (Hint: Add one new variable $y$ and replace each clause $(\ell_0 \vee \ell_1 \vee \ell_2)$ by $(\ell_0 \vee \ell_1 \vee \ell_2 \vee y)$.)
        \item \FlipFourSAT\ can be reduced to \FlipThreeSAT\ by the same method we used to reduce \SAT\ to \ThreeSAT, and thus \FlipThreeSAT\ is also $\NPsearch$-complete.  Using this fact,
        prove that \GraphThreeColoring\ is $\NPsearch$-complete. (Hint: use the same construction as we used in the reduction from \ThreeSAT\ to \IndependentSet, except add one extra vertex that's connected to all of the variable-gadget-vertices.)
        \item (optional\footnote{This problem won't make a difference between N, L, R-, and R grades. As this problem is purely extra credit, course staff will deprioritize questions about this problem at office hours and on Ed.}) Fill in the omitted details of the reduction from \FlipFourSAT\ to \FlipThreeSAT, with its proof of correctness.
\end{enumerate}

\item (Factoring if $\NPsearch\subset \Psearch$) Every time your browser or one of your apps makes a secure connection to a website or server, it uses cryptography to make sure that eavesdroppers to the network cannot decipher the information you are transmitting. Factoring a large number is an important part of cryptography, and the basis for the RSA encryption. Lets define the computational problem formally.
\compprob{\Factoring}
        {A natural number $N$ in base 2}
        {Two natural numbers $P,Q$ in base 2, such that $PQ=N$ (if exists)}

Show that if $\NPsearch\subset \Psearch$, then there is an algorithm to solve \Factoring in time $O((\log N)^c)$, for some $c=O(1)$. Be careful about the length of input in base 2.

\textit{(Fun Fact: A remarkable property of quantum computers is that they can solve \Factoring in polynomial time, via Shor's algorithm. This does not mean that quantum computers can solve any problem in $\NPsearch$ in polynomial time, as \Factoring is not believed to be $\NPsearch$-complete).
}

\item (reflection) Describe one theoretical idea from this course that you have found beautiful, and explain why it is beautiful to you.  Your answer should: (1) explain the idea in a way that could be understood by a classmate who has taken classes CS20 and CS50 but has not yet taken this class and (2) address how this beauty is similar to or different from other kinds of beauty that human beings encounter.

Quick note on grading: Good responses are usually about a paragraph, with something like 7 or 8 sentences. Most importantly, please make sure your answer is specific to this class and your experiences in it. If your answer could have been edited lightly to apply to another class at Harvard, points will be taken off.

\textit{Note: As with the previous psets, you may include your answer in your PDF submission, but the answer should ultimately go into a separate Gradescope submission form.}

\item Once you're done with this problem set, please fill out \href{https://forms.gle/UWfAWiH9SdEjMQos7}{this survey} so that we can gather students' thoughts on the problem set, and the class in general. It's not required, but we really appreciate all responses!

\end{enumerate}
\end{document}